\chapter{КОДЫН ЖИШЭЭ БА ӨГӨГДЛИЙН СХЕМ}
\begin{sloppy}

Энэхүү хавсралтад дипломын ажлын үндсэн логикийг илэрхийлэх representative кодын жишээ болон өгөгдлийн сангийн схемийг оруулав. Доорх листингүүд нь байршлын шалгалт, AR placement, авдрын interaction, reward claim хамгаалалт, өгөгдлийн бүтэц зэрэг системийн цөм хэсгүүдийг тайлбарлах зорилготой.

\section{Ойролцоох байршил шалгах код}

\begin{lstlisting}[language={[Sharp]C}, caption=TourismSpotManager - proximity шалгах үндсэн логик]
private void UpdateNearbySpots()
{
    ResetRuntimeState();

    if (locationTracker == null || !locationTracker.IsLocationReady || tourismSpots.Count == 0)
    {
        return;
    }

    foreach (TourismSpot spot in tourismSpots)
    {
        float distanceMeters = CalculateDistanceMeters(
            locationTracker.Latitude,
            locationTracker.Longitude,
            spot.latitude,
            spot.longitude);

        if (distanceMeters <= spot.radiusMeters)
        {
            isInsideAnySpot = true;
            currentNearbySpot = spot;
        }
    }
}
\end{lstlisting}

\section{AR орчинд авдар байрлуулах код}

\begin{lstlisting}[language={[Sharp]C}, caption=ARChestSpawner - fallback placement логик]
private void SpawnFallbackChest()
{
    Vector3 fallbackPosition = GetFallbackChestPosition();
    Quaternion rotation = GetChestFacingRotation();

    GameObject spawnedChest = Instantiate(chestPrefab, fallbackPosition, rotation);
    PrepareSpawnedChest(spawnedChest);

    hasSpawned = true;
}
\end{lstlisting}

\section{Авдар нээх ба reward гаргах код}

\begin{lstlisting}[language={[Sharp]C}, caption=ARChestInteraction - chest open хамгаалалт]
private void OpenChest()
{
    if (isOpening || isOpened)
    {
        return;
    }

    isOpening = true;
    isOpened = true;
    DisableChestColliders();
    PlayOpenAnimation();

    if (rewardRevealRoutine != null)
    {
        StopCoroutine(rewardRevealRoutine);
    }

    rewardRevealRoutine = StartCoroutine(RevealRewardCollectible());
}
\end{lstlisting}

\section{Reward claim endpoint-ийн жишээ}

\begin{lstlisting}[language=JavaScript, caption=Express endpoint - давхардал хамгаалалттай reward claim]
app.post("/spots/:id/claim", async (req, res) => {
  const userId = req.user.id;
  const tourismSpotId = req.params.id;

  const existing = await prisma.userReward.findUnique({
    where: {
      userId_tourismSpotId: { userId, tourismSpotId }
    }
  });

  if (existing) {
    return res.json({ status: "alreadyClaimed", rewardId: existing.rewardId });
  }

  const spot = await prisma.tourismSpot.findUnique({
    where: { id: tourismSpotId }
  });

  const claim = await prisma.userReward.create({
    data: {
      userId,
      tourismSpotId,
      rewardId: spot.rewardId
    }
  });

  return res.json({ status: "claimed", claim });
});
\end{lstlisting}

\section{Өгөгдлийн сангийн зохиомж ба хүснэгтүүд}
\begin{sloppy}

Системийн өгөгдлийн сан нь `User`, `TourismSpot`, `Reward`, `UserReward` гэсэн үндсэн entity-үүдэд тулгуурлана. Энэхүү загварыг PostgreSQL реляцийн мэдээллийн сан дээр Prisma ORM ашиглан хэрэгжүүлсэн бөгөөд reward collection-ийн бизнес дүрмийг өгөгдлийн түвшинд баталгаажуулахад чиглэсэн.

\subsection{Хүснэгтүүдийн бүтэц}

\begin{table}[H]
    \centering
    \caption{Users (Хэрэглэгчид) хүснэгтийн бүтэц}
    \label{tab:users_table}
    \begin{tabular}{|l|l|l|}
        \hline
        \textbf{Талбарын нэр} & \textbf{Төрөл} & \textbf{Тайлбар} \\ \hline
        id & String (CUID) & Хэрэглэгчийн цор ганц ID \\ \hline
        deviceId & String (Unique) & Төхөөрөмжийн танигч \\ \hline
        displayName & String & Хэрэглэгчийн харагдах нэр \\ \hline
        lastLoginAt & DateTime & Хамгийн сүүлд session сэргээсэн огноо \\ \hline
    \end{tabular}
\end{table}

\begin{table}[H]
    \centering
    \caption{TourismSpots (Аялал жуулчлалын байршил) хүснэгтийн бүтэц}
    \label{tab:tourism_spots_table}
    \begin{tabular}{|l|l|l|}
        \hline
        \textbf{Талбарын нэр} & \textbf{Төрөл} & \textbf{Тайлбар} \\ \hline
        id & String (CUID) & Байршлын цор ганц ID \\ \hline
        name & String & Байршлын нэр \\ \hline
        latitude & Decimal & Өргөрөг \\ \hline
        longitude & Decimal & Уртраг \\ \hline
        radiusMeters & Float & Идэвхжих радиус \\ \hline
        rewardId & String & Холбогдох шагналын ID \\ \hline
    \end{tabular}
\end{table}

\begin{table}[H]
    \centering
    \caption{Rewards (Шагнал) хүснэгтийн бүтэц}
    \label{tab:rewards_table}
    \begin{tabular}{|l|l|l|}
        \hline
        \textbf{Талбарын нэр} & \textbf{Төрөл} & \textbf{Тайлбар} \\ \hline
        id & String (CUID) & Шагналын цор ганц ID \\ \hline
        name & String & Шагналын нэр \\ \hline
        description & String & Соёлын болон түүхэн тайлбар \\ \hline
        imageUrl & String & Reward дүрслэлийн URL \\ \hline
        previewPrefabKey & String & Unity prefab-ийн түлхүүр \\ \hline
    \end{tabular}
\end{table}

\begin{table}[H]
    \centering
    \caption{UserRewards (Хэрэглэгчийн цуглуулга) хүснэгтийн бүтэц}
    \label{tab:user_rewards_table}
    \begin{tabular}{|l|l|l|}
        \hline
        \textbf{Талбарын нэр} & \textbf{Төрөл} & \textbf{Тайлбар} \\ \hline
        id & String (CUID) & Бүртгэлийн цор ганц ID \\ \hline
        userId & String & Хэрэглэгчийн ID \\ \hline
        rewardId & String & Авсан reward-ийн ID \\ \hline
        tourismSpotId & String & Reward авсан байршлын ID \\ \hline
        claimedAt & DateTime & Reward claim хийсэн огноо \\ \hline
    \end{tabular}
\end{table}

\subsection{Хүснэгт хоорондын хамаарал}

`TourismSpot` бүр нэг `Reward`-тай холбогдоно. `UserReward` хүснэгт нь `User`, `TourismSpot`, `Reward` гурвыг холбож, аль хэрэглэгч аль байршлаас аль reward авсныг хадгална. `@@unique([userId, tourismSpotId])` constraint нь нэг хэрэглэгч нэг байршлын reward-ийг давхар авахыг хориглодог. Ингэснээр reward collection-ийн бизнес дүрэм зөвхөн клиент талын интерфэйс дээр бус, өгөгдлийн сангийн түвшинд ч хамгаалагдана.

\end{sloppy}


\end{sloppy}
